% Options for packages loaded elsewhere
\PassOptionsToPackage{unicode}{hyperref}
\PassOptionsToPackage{hyphens}{url}
\documentclass[
]{article}
\usepackage{xcolor}
\usepackage{amsmath,amssymb}
\setcounter{secnumdepth}{-\maxdimen} % remove section numbering
\usepackage{iftex}
\ifPDFTeX
  \usepackage[T1]{fontenc}
  \usepackage[utf8]{inputenc}
  \usepackage{textcomp} % provide euro and other symbols
\else % if luatex or xetex
  \usepackage{unicode-math} % this also loads fontspec
  \defaultfontfeatures{Scale=MatchLowercase}
  \defaultfontfeatures[\rmfamily]{Ligatures=TeX,Scale=1}
\fi
\usepackage{lmodern}
\ifPDFTeX\else
  % xetex/luatex font selection
\fi
% Use upquote if available, for straight quotes in verbatim environments
\IfFileExists{upquote.sty}{\usepackage{upquote}}{}
\IfFileExists{microtype.sty}{% use microtype if available
  \usepackage[]{microtype}
  \UseMicrotypeSet[protrusion]{basicmath} % disable protrusion for tt fonts
}{}
\makeatletter
\@ifundefined{KOMAClassName}{% if non-KOMA class
  \IfFileExists{parskip.sty}{%
    \usepackage{parskip}
  }{% else
    \setlength{\parindent}{0pt}
    \setlength{\parskip}{6pt plus 2pt minus 1pt}}
}{% if KOMA class
  \KOMAoptions{parskip=half}}
\makeatother
\usepackage{longtable,booktabs,array}
\newcounter{none} % for unnumbered tables
\usepackage{calc} % for calculating minipage widths
% Correct order of tables after \paragraph or \subparagraph
\usepackage{etoolbox}
\makeatletter
\patchcmd\longtable{\par}{\if@noskipsec\mbox{}\fi\par}{}{}
\makeatother
% Allow footnotes in longtable head/foot
\IfFileExists{footnotehyper.sty}{\usepackage{footnotehyper}}{\usepackage{footnote}}
\makesavenoteenv{longtable}
\setlength{\emergencystretch}{3em} % prevent overfull lines
\providecommand{\tightlist}{%
  \setlength{\itemsep}{0pt}\setlength{\parskip}{0pt}}
% ===== unicode + compile header for ZIMMERMAN_LAW_OF_GRAVITY =====
% TO COMPILE ON OVERLEAF: upload this .tex + a figures/ folder containing
%   fig1_rar.png  fig2_a0z.png  fig3_btfr.png
% Compile with XeLaTeX (recommended) or pdfLaTeX. Download the PDF and publish
% manually on Zenodo. (Not test-compiled locally; unicode is mapped below.)
\usepackage{amssymb}
\usepackage{newunicodechar}
\providecommand{\textsubscript}[1]{\ensuremath{_{\mbox{\scriptsize #1}}}}
\providecommand{\textonehalf}{\ensuremath{\frac12}}
\newunicodechar{₀}{\ensuremath{_{0}}}
\newunicodechar{₁}{\ensuremath{_{1}}}
\newunicodechar{₃}{\ensuremath{_{3}}}
\newunicodechar{₆}{\ensuremath{_{6}}}
\newunicodechar{ₐ}{\ensuremath{_{a}}}
\newunicodechar{¹}{\ensuremath{^{1}}}
\newunicodechar{²}{\ensuremath{^{2}}}
\newunicodechar{⁰}{\ensuremath{^{0}}}
\newunicodechar{⁴}{\ensuremath{^{4}}}
\newunicodechar{⁵}{\ensuremath{^{5}}}
\newunicodechar{⁻}{\ensuremath{^{-}}}
\newunicodechar{Λ}{\ensuremath{\Lambda}}
\newunicodechar{ρ}{\ensuremath{\rho}}
\newunicodechar{π}{\ensuremath{\pi}}
\newunicodechar{σ}{\ensuremath{\sigma}}
\newunicodechar{Υ}{\ensuremath{\Upsilon}}
\newunicodechar{χ}{\ensuremath{\chi}}
\newunicodechar{Ω}{\ensuremath{\Omega}}
\newunicodechar{Δ}{\ensuremath{\Delta}}
\newunicodechar{γ}{\ensuremath{\gamma}}
\newunicodechar{α}{\ensuremath{\alpha}}
\newunicodechar{µ}{\ensuremath{\mu}}
\newunicodechar{√}{\ensuremath{\surd}}
\newunicodechar{×}{\ensuremath{\times}}
\newunicodechar{−}{\ensuremath{-}}
\newunicodechar{·}{\ensuremath{\cdot}}
\newunicodechar{∝}{\ensuremath{\propto}}
\newunicodechar{≈}{\ensuremath{\approx}}
\newunicodechar{≳}{\ensuremath{\gtrsim}}
\newunicodechar{≥}{\ensuremath{\geq}}
\newunicodechar{⇒}{\ensuremath{\Rightarrow}}
\newunicodechar{→}{\ensuremath{\rightarrow}}
\newunicodechar{↔}{\ensuremath{\leftrightarrow}}
\newunicodechar{⟺}{\ensuremath{\Longleftrightarrow}}
\newunicodechar{⋆}{\ensuremath{\star}}
\newunicodechar{★}{\ensuremath{\star}}
\newunicodechar{✓}{\ensuremath{\checkmark}}
\newunicodechar{✗}{\ensuremath{\times}}
\newunicodechar{◐}{\ensuremath{\circ}}
\newunicodechar{½}{\ensuremath{\frac{1}{2}}}
\newunicodechar{§}{\S{}}
\newunicodechar{Ö}{\"{O}}
\newunicodechar{Ü}{\"{U}}
\newunicodechar{é}{\'{e}}
\newunicodechar{³}{\ensuremath{^{3}}}
\newunicodechar{ä}{\"{a}}
\newunicodechar{Ȳ}{\ensuremath{\bar{Y}}}
\newunicodechar{δ}{\ensuremath{\delta}}
\newunicodechar{θ}{\ensuremath{\theta}}
\newunicodechar{κ}{\ensuremath{\kappa}}
\newunicodechar{λ}{\ensuremath{\lambda}}
\newunicodechar{μ}{\ensuremath{\mu}}
\newunicodechar{φ}{\ensuremath{\varphi}}
\newunicodechar{–}{\textendash{}}
\newunicodechar{—}{\textemdash{}}
\newunicodechar{…}{\ldots{}}
\newunicodechar{′}{'}
\newunicodechar{⁶}{\ensuremath{^{6}}}
\newunicodechar{⁷}{\ensuremath{^{7}}}
\newunicodechar{₂}{\ensuremath{_{2}}}
\newunicodechar{₄}{\ensuremath{_{4}}}
\newunicodechar{ℏ}{\ensuremath{\hbar}}
\newunicodechar{∇}{\ensuremath{\nabla}}
\newunicodechar{≔}{\ensuremath{:=}}
\newunicodechar{≤}{\ensuremath{\leq}}
\newunicodechar{≲}{\ensuremath{\lesssim}}
\newunicodechar{⊙}{\ensuremath{\odot}}
\newunicodechar{Σ}{\ensuremath{\Sigma}}
\newunicodechar{η}{\ensuremath{\eta}}
\newunicodechar{ν}{\ensuremath{\nu}}
\newunicodechar{″}{''}
\newunicodechar{⇔}{\ensuremath{\Leftrightarrow}}
\newunicodechar{±}{\ensuremath{\pm}}
\newunicodechar{¼}{\ensuremath{\tfrac{1}{4}}}
\newunicodechar{β}{\ensuremath{\beta}}
\newunicodechar{∈}{\ensuremath{\in}}
\newunicodechar{∞}{\ensuremath{\infty}}
\newunicodechar{≠}{\ensuremath{\neq}}
% --- v3 erratum build: additional symbol coverage ---
\newunicodechar{ω}{\ensuremath{\omega}}
\newunicodechar{τ}{\ensuremath{\tau}}
\newunicodechar{∫}{\ensuremath{\int}}
\newunicodechar{∂}{\ensuremath{\partial}}
\newunicodechar{≫}{\ensuremath{\gg}}
\newunicodechar{⟨}{\ensuremath{\langle}}
\newunicodechar{⟩}{\ensuremath{\rangle}}
\newunicodechar{⁸}{\ensuremath{^{8}}}
\usepackage{bookmark}
\IfFileExists{xurl.sty}{\usepackage{xurl}}{} % add URL line breaks if available
\urlstyle{same}
\hypersetup{
  hidelinks,
  pdfcreator={LaTeX via pandoc}}

\author{}
\date{}

\begin{document}

\section{\texorpdfstring{Does the Galaxy Acceleration Scale Track Cosmic
Dark Energy? A Non-Circular Cross-Scale \(a_0(z)\) Test of de
Sitter--Unruh Modified Inertia --- the \(z=0\) Tie, the High-\(z\) BTFR
Forecast, and Why \(M_\mathrm{bar}\) Calibration (Not Rotator Count) Is
Decisive}{Does the Galaxy Acceleration Scale Track Cosmic Dark Energy? A Non-Circular Cross-Scale a\_0(z) Test of de Sitter--Unruh Modified Inertia --- the z=0 Tie, the High-z BTFR Forecast, and Why M\_\{\textbackslash rm bar\} Calibration (Not Rotator Count) Is Decisive}}\label{does-the-galaxy-acceleration-scale-track-cosmic-dark-energy-a-non-circular-cross-scale-a_0z-test-of-de-sitterunruh-modified-inertia-the-z0-tie-the-high-z-btfr-forecast-and-why-m_rm-bar-calibration-not-rotator-count-is-decisive}

\textbf{Carl Zimmerman} Briar Creek Tech · carl@briarcreektech.com

\begin{center}\rule{0.5\linewidth}{0.5pt}\end{center}

\subsection{Abstract}\label{abstract}

In a de Sitter--Unruh modified-inertia (MI) framework the
low-acceleration scale \(a_0\) is fixed to the cosmological horizon by
\(a_0=cH_\Lambda/Z\) with \(Z=\sqrt{32\pi/3}=5.78881\), giving the
pure--dark-energy-density (canonical) footing
\(a_0^\mathrm{canon}=c^2\sqrt{\Lambda/32\pi}=9.355\times10^{-11}\,{\rm m\,s^{-2}}\)
or the total-density (ALT) footing
\(a_0^\mathrm{ALT}=cH_0/Z=1.1305\times10^{-10}\,{\rm m\,s^{-2}}\), two
footings \(20.9\%\) apart that are carried on every number here. The
framework does \textbf{not} rewrite Type Ia supernova standardization
and gives \textbf{no} new distance formula --- its cosmological
background is General Relativity. Its distinctive content is that the
dark-energy term in the Friedmann equation is not a free fit but the
\emph{galaxy} acceleration scale:
\(H^2(z)=\tfrac{8\pi G}{3}\rho_m(z)+Z^2a_0^2(z)/c^2\), with \(a_0(z)\)
measurable from galaxy dynamics and \(Z\) a single posited constant.
This ties a galaxy-kinematic quantity to a cosmic-distance quantity
through \(a_0(z)=\tfrac{c}{2}\sqrt{G\rho_\mathrm{DE}(z)}\) (equivalently
\(a_0=c^2\sqrt{\Lambda_\mathrm{eff}/32\pi}\) with
\(\Lambda_\mathrm{eff}=8\pi G\rho_\mathrm{DE}/c^2\)),
i.e.~\(a_0(z)/a_0(0)=\sqrt{\rho_\mathrm{DE}(z)/\rho_\mathrm{DE,0}}\) --- a
bridge \(\Lambda\)CDM structurally lacks. It adds no dark-energy
parameters beyond the GR background: the same \(\rho_\mathrm{DE}(z)\) that
fits cosmic distances must fit the galaxy \(a_0(z)\). The equality
\(a_0\!\leftrightarrow\!\Lambda\) at fixed \(z\) is a definition and is
circular; the \emph{non-circular} test is the joint \(z\)-tracking of
two independent datasets. We report a method + forecast
\textbf{non-detection with a live cosmology-side hint}. (i) The \(z=0\)
tie holds: \(\Lambda\)-blind SPARC
\(a_0(0)=1.181\times10^{-10}\,(\pm16\%)\) versus cosmic
\((c/2)\sqrt{G\rho_\mathrm{DE,0}}\) gives ratio \(1.26\) (\(1.44\sigma\),
Planck \(H_0\)) or \(1.17\) (\(0.95\sigma\), SH0ES \(H_0\)). (ii) The
DESI DR2 \(w_0w_a\) posterior propagated through the framework relation
yields \(a_0(z{=}3)/a_0(0)\simeq0.74\) (Pantheon+ \(0.775\), DESY5
\(0.737\), Union3 \(0.707\); a \(\sim2.2\sigma\) decline) --- a
consequence of the framework relation plus the DESI \(w(z)\), not an
independent detection. (iii) On the galaxy side the distinctive
\(0.60\)--\(0.75\) decline is neither detected nor excluded:
separability \(S\simeq0.8\sigma\). One clean deep-MOND \(z=3.25\)
rotator (the Big Wheel) gives \(a_0(3.25)/a_0(0)=1.31\,(+0.93/-0.52)\),
consistent with a constant and excluding the ALT-\(cH_0\)
\(\sim5\times\) \emph{rise} at \(\sim2\sigma\), but unable to separate
the \(0.70\) decline from flat. The central finding is methodological:
in the baryonic Tully--Fisher (BTFR) estimator \(a_0\) is algebraically
\(1{:}1\) degenerate with baryonic mass
(\(\Delta\log a_0=-\Delta\log M_\mathrm{bar}\) at fixed \(V\)), so the
dominant \(M_\mathrm{bar}\) systematic is common-mode and does \textbf{not}
beat down with rotator count: with a \(\gtrsim0.08\)-dex correlated
\(M_\mathrm{bar}\) floor the test never reaches \(3\sigma\) at any \(N\).
The decisive lever is \(M_\mathrm{bar}\) calibration to \(\lesssim0.05\)
dex common-mode, not the number of galaxies, and the \(a_0\)-separable
observable that breaks the degeneracy is the full rotation-curve /
radial-acceleration transition, not the BTFR zero-point. The value of
\(a_0\) and the constant \(Z\) are posited, not derived; the \(a_0\)
magnitude inherits \(Z\), so only the ratio/tracking is under test.

\begin{center}\rule{0.5\linewidth}{0.5pt}\end{center}

\subsection{1. Introduction}\label{introduction}

\subsubsection{1.1 The framework and the horizon
scale}\label{the-framework-and-the-horizon-scale}

Modified inertia proposes that the inertial response of a body to a
force weakens below a critical acceleration \(a_0\), rather than the
gravitational field being modified. In the de Sitter--Unruh realization,
an accelerated observer in a universe with a cosmological horizon sees a
floor in the vacuum's response set by the horizon temperature, and the
transition scale is tied to the horizon:

\[a_0=\frac{cH_\Lambda}{Z}=\frac{c^2\sqrt{\Lambda/3}}{Z},\qquad Z=\sqrt{\tfrac{32\pi}{3}}=5.78881,\]

with \(H_\Lambda\equiv c\sqrt{\Lambda/3}\) the de Sitter (pure
dark-energy) Hubble rate. The Planck 2018 value
\(\Lambda=1.089\times10^{-52}\,{\rm m^{-2}}\) gives the
\textbf{canonical footing}
\(a_0^\mathrm{canon}=c^2\sqrt{\Lambda/32\pi}=9.355\times10^{-11}\,{\rm m\,s^{-2}}\)
(pure \(\rho_\mathrm{DE}\)). A second \textbf{ALT footing} replaces the
pure-\(\Lambda\) density by the total present-day critical density,
\(a_0^\mathrm{ALT}=cH_0/Z=1.1305\times10^{-10}\,{\rm m\,s^{-2}}\). The two
differ by \(20.9\%\); both are carried on every dimensional number
below. Nothing here derives the value of \(a_0\), the constant \(Z\), or
the sign of the inertial correction: those remain postulates. What the
framework ties to the horizon is the \emph{scale} of \(a_0\).

\subsubsection{\texorpdfstring{1.2 The parameter-free replacement of the
free-\(w\)
fit}{1.2 The parameter-free replacement of the free-w fit}}\label{the-parameter-free-replacement-of-the-free-w-fit}

The framework keeps General Relativity for the cosmological background
and therefore has no native supernova formula: the light-curve \(\to\)
distance standardization is unchanged, and the Friedmann equation is the
standard one. Its single distinctive move is that the dark-energy term
is not a free fit but the galaxy acceleration scale,

\[H^2(z)=\frac{8\pi G}{3}\,\rho_m(z)+\frac{Z^2a_0^2(z)}{c^2},\qquad \frac{Z^2a_0^2}{c^2}=H_\Lambda^2=\frac{\Lambda c^2}{3}\;\;(z=0),\]

so that, inverting \(Z^2a_0^2/c^2=\tfrac{8\pi G}{3}\rho_\mathrm{DE}\) with
\(Z^2=32\pi/3\), the dark-energy mass density is
\(\rho_\mathrm{DE}(z)=4a_0^2(z)/(Gc^2)\) and

\[\boxed{\;a_0(z)=\tfrac{c}{2}\sqrt{G\,\rho_\mathrm{DE}(z)}\;=\;c^2\sqrt{\Lambda_\mathrm{eff}(z)/32\pi}\quad\Longrightarrow\quad \frac{a_0(z)}{a_0(0)}=\sqrt{\frac{\rho_\mathrm{DE}(z)}{\rho_\mathrm{DE,0}}}\;}\]

The SI form \(\tfrac{c}{2}\sqrt{G\rho_\mathrm{DE}}\) uses the dark-energy
\emph{mass} density \(\rho_\mathrm{DE}\) (kg m\(^{-3}\)); the equivalent
geometric form uses \(\Lambda_\mathrm{eff}\equiv8\pi G\rho_\mathrm{DE}/c^2\)
(units m\(^{-2}\)). Both give \(a_0^\mathrm{canon}=9.355\times10^{-11}\) at
\(z=0\) on Planck \(\rho_\mathrm{DE,0}\); the two conventions must not be
mixed (\(c^2\sqrt{\rho_\mathrm{DE}/32\pi}\) with a \emph{mass} density is
dimensionally inconsistent). The redshift-tracking ratio
\(a_0(z)/a_0(0)=\sqrt{\rho_\mathrm{DE}(z)/\rho_\mathrm{DE,0}}\) is
convention-independent.

On the supernova side this adds \textbf{no free dark-energy parameter}
--- against \(\Lambda\)CDM's \(\Omega_\Lambda\) (and, in the evolving-DE
extension, \(w_0,w_a\)). The framework does not, however, produce a
parameter-free dark-energy history from nothing: it \emph{consumes}
whatever \(\rho_\mathrm{DE}(z)\) the cosmic data prefer (itself the
\(w_0,w_a\) fit to distances) and requires that the \emph{same} function
reproduce the galaxy-side \(a_0(z)\). The content is therefore a
cross-\emph{dataset} constraint --- one \(\rho_\mathrm{DE}(z)\) must
simultaneously fit cosmic distances and galaxy kinematics --- not a
first-principles prediction of \(w(z)\). At \(z=0\) the framework is
identical to \(\Lambda\)CDM by construction
(\(Z^2a_0^2/c^2=\Lambda c^2/3\)). It differs from \(\Lambda\)CDM only if
\(a_0(z)\) evolves --- that is, only if the galaxy acceleration scale
changes with redshift in lock-step with the cosmic dark-energy density.

\subsubsection{1.3 The non-circular cross-scale
idea}\label{the-non-circular-cross-scale-idea}

\(\Lambda\)CDM gives no reason for a galaxy's internal acceleration
scale to track cosmic expansion; the framework forces exactly that link.
Crucially, the equality of a galaxy \(a_0\) and a cosmic \(\Lambda\) at
a \emph{single} epoch is a definition --- one can always convert one
into the other, and reading agreement there is circular. The genuine,
non-circular test is whether \textbf{two independently measured
functions of redshift} --- \(a_0\) from galaxy kinematics and
\(\rho_\mathrm{DE}\) from cosmic distances --- track each other as \(z\)
varies. That \(z\)-tracking is a prediction \(\Lambda\)CDM cannot make
and the framework cannot avoid. The framework \emph{inherits} the
dark-energy history \(w(z)\) (it does not predict \(w(z)\)); its content
is the tie between the two datasets, so any decline it exhibits is a
consequence of the framework relation combined with the measured
\(w(z)\), not an independent statement about dark energy.

\subsubsection{1.4 Credit for the kernel, and the status of this
paper}\label{credit-for-the-kernel-and-the-status-of-this-paper}

The low-acceleration interpolation used to read \(a_0\) from rotation
curves, \(\nu(y)=\sqrt{1+1/y}\) with \(y=g_\mathrm{bar}/a_0\), is Milgrom's
(1999, \emph{Phys. Lett. A} \textbf{253}, 273, Eq. 9), in the original
MOND context of Milgrom (1983). Milgrom's coefficient was
\(2cH_\Lambda\) rather than \(cH_\Lambda/Z\); the deep-MOND limit
\(V^4=a_0\,GM_\mathrm{bar}\) that anchors the galaxy-side measurement here
is the baryonic Tully--Fisher relation (McGaugh 2005, 2012; Lelli,
McGaugh \& Schombert 2016). The distinctive content the present
framework contributes is the \(cH_\Lambda/Z\) coefficient and the
\(a_0\!\leftrightarrow\!\rho_\mathrm{DE}(z)\) tie under test. This is a
\textbf{method + forecast paper reporting a non-detection} with a live
cosmology-side hint and a strategic finding; it is not, and does not
claim to be, a detection. All load-bearing numbers are produced by
committed, exit-0 scripts (Appendix A), and both footings are carried
throughout.

\begin{center}\rule{0.5\linewidth}{0.5pt}\end{center}

\subsection{2. The relation and why it is
non-circular}\label{the-relation-and-why-it-is-non-circular}

\textbf{Terminology.} \emph{Footing} --- the choice of density that sets
the \(a_0\) \emph{magnitude}: canonical (pure \(\rho_\mathrm{DE}\),
\(9.36\times10^{-11}\)) or ALT (total present-day density,
\(1.13\times10^{-10}\)); the two are \(20.9\%\) apart and cancel in any
ratio. \emph{Deep-MOND} --- the low-acceleration regime \(g\ll a_0\),
where \(\nu\to\sqrt{a_0/g}\) and \(V^4=a_0GM_\mathrm{bar}\) holds with
fitted slope exactly \(4\). \emph{BTFR zero-point} --- the normalization
of \(V^4=a_0GM_\mathrm{bar}\); in the deep-MOND limit that zero-point
\emph{is} \(a_0\), so reading \(a_0(z)\) and reading the BTFR zero-point
are the same operation. \emph{Common-mode} --- a systematic applied
identically across a sample (one \(\alpha_\mathrm{CO}\)/IMF/dust
prescription), which does not average down with sample size \(N\).

\subsubsection{2.1 The two-dataset bridge}\label{the-two-dataset-bridge}

Write the framework tie as a bridge between measurements made on wholly
different physical scales:

\[\underbrace{a_0(z)}_{\text{galaxy kinematics: BTFR zero-point}}\;=\;\tfrac{c}{2}\sqrt{G\,\rho_\mathrm{DE}(z)}\;=\;\underbrace{\tfrac{c}{2}\sqrt{G\,\rho_\mathrm{DE}(z)}}_{\text{cosmic distances: SNe / DESI }\rho_\mathrm{DE}(z)}.\]

The left side is read from the internal dynamics of individual galaxies
(parsec--kiloparsec scales); the right side is read from the expansion
history (gigaparsec scales). \(\Lambda\)CDM contains no relation forcing
these to agree at all, let alone to co-evolve. The framework forces both
their \(z=0\) equality and their joint \(z\)-dependence.

\subsubsection{2.2 Definition versus test}\label{definition-versus-test}

The \(z=0\) equality \(a_0(0)=(c/2)\sqrt{G\rho_\mathrm{DE,0}}\) is
definitional: given a cosmic \(\rho_\mathrm{DE,0}\) one \emph{defines} the
horizon \(a_0\), so a \(z=0\) match is a consistency knot, not an
independent success. Its value is as the anchor the bridge extends from
--- if it failed badly the whole construction would be dead on arrival
--- but it carries no evidential weight beyond internal consistency. The
non-circular content is entirely in the \emph{slope}: does
galaxy-measured \(a_0(z)\) decline (or stay flat, or rise) in the same
way the independently measured cosmic \(\rho_\mathrm{DE}(z)\) does? Because
the two datasets never share inputs, a wrong background assumption
shifts both together and cannot manufacture a spurious co-evolution
(quantified in §4.4).

\subsubsection{2.3 What the framework predicts, and what it
inherits}\label{what-the-framework-predicts-and-what-it-inherits}

At fixed \(z\) the framework prediction is fully specified once
\(\rho_\mathrm{DE}(z)\) is chosen. Two reference cosmic tracks bracket the
question:

{\def\LTcaptype{none} % do not increment counter
\begin{longtable}[]{@{}
  >{\raggedright\arraybackslash}p{(\linewidth - 6\tabcolsep) * \real{0.2500}}
  >{\raggedright\arraybackslash}p{(\linewidth - 6\tabcolsep) * \real{0.2500}}
  >{\raggedright\arraybackslash}p{(\linewidth - 6\tabcolsep) * \real{0.2500}}
  >{\raggedright\arraybackslash}p{(\linewidth - 6\tabcolsep) * \real{0.2500}}@{}}
\toprule\noalign{}
\begin{minipage}[b]{\linewidth}\raggedright
track
\end{minipage} & \begin{minipage}[b]{\linewidth}\raggedright
\((w_0,w_a)\)
\end{minipage} & \begin{minipage}[b]{\linewidth}\raggedright
\(a_0(3)/a_0(0)\)
\end{minipage} & \begin{minipage}[b]{\linewidth}\raggedright
\(a_0(3.25)/a_0(0)\)
\end{minipage} \\
\midrule\noalign{}
\endhead
\bottomrule\noalign{}
\endlastfoot
flat-\(\Lambda\) (\(w=-1\)) & \((-1.00,\,0.00)\) & \(1.000\) (flat) &
\(1.000\) \\
DESI-DR2 evolving & \((-0.75,\,-0.90)\) & \(0.712\) (declining) &
\(0.685\) \\
\emph{(ALT \(cH(z)\) evolution model, for contrast)} &
\(a_0\propto cH(z)=cH_0E(z)\) & \(4.57\) (rising) & \(4.99\) \\
\end{longtable}
}

with CPL history
\(\rho_\mathrm{DE}(z)/\rho_\mathrm{DE,0}=(1+z)^{3(1+w_0+w_a)}\exp[-3w_az/(1+z)]\).
The framework's distinctive canonical-footing prediction is a
\textbf{mild decline}, \(a_0(3)/a_0(0)\approx0.60\)--\(0.75\), if the
DESI-preferred evolving dark energy is real. If instead \(w\to-1\)
(flat), the framework reduces to \(\Lambda\)CDM at all \(z\) and the
distinctive signal vanishes, leaving only the \(z=0\) tie. The third row
is \textbf{not} the ALT \(z=0\) normalization
(\(a_0^\mathrm{ALT}=1.131\times10^{-10}\)) but a distinct \emph{evolution
model} in which \(a_0\) tracks the total-density Hubble rate \(cH(z)\)
and therefore \emph{rises} steeply --- the opposite sign --- and it is
the one branch the current galaxy data can already address (§4.3). The
ALT \(z=0\) normalization itself, if it instead tracked
\(\rho_\mathrm{DE}(z)\), would decline like the canonical footing and is
\emph{not} excluded; only the rising \(cH(z)\) evolution is. The
flat--DESI separation is only \(0.147\) dex at \(z=3\) (\(0.164\) dex at
\(z=3.25\)); that small gap is what the galaxy data must resolve.

A theory-side caveat bounds the whole construction. The horizon relation
\(a_0=cH_\Lambda/Z\) is derived for a \emph{static} de Sitter horizon
(constant \(\Lambda\)). Promoting it to
\(a_0(z)=\tfrac{c}{2}\sqrt{G\rho_\mathrm{DE}(z)}\) for an evolving equation
of state (\(w\ne-1\)) is an \textbf{adiabatic / instantaneous-horizon
ansatz} --- the assumption that the inertia scale tracks the
instantaneous dark-energy density --- not a result derived from the de
Sitter--Unruh construction, whose horizon is that of a
constant-\(\Lambda\) background. The \(z\)-tracking under test is
therefore a test of this ansatz combined with the measured \(w(z)\); a
null would disfavour the ansatz, not the static-horizon relation it
reduces to at \(z=0\).

\begin{center}\rule{0.5\linewidth}{0.5pt}\end{center}

\subsection{\texorpdfstring{3. The \(z=0\) tie and the cosmology
side}{3. The z=0 tie and the cosmology side}}\label{the-z0-tie-and-the-cosmology-side}

\subsubsection{\texorpdfstring{3.1 The \(z=0\)
tie}{3.1 The z=0 tie}}\label{the-z0-tie}

The \(\Lambda\)-blind galaxy \(a_0(0)\) is taken from the gas-dominated
SPARC BTFR zero-point measured without any cosmological-\(\Lambda\)
input (the companion \(a_0\)-line/\(\Lambda\) paper, Zimmerman 2026, DOI
10.5281/zenodo.21419735):
\(a_0(0)=1.181\times10^{-10}\,{\rm m\,s^{-2}}\,(\pm16\%)\). The
cosmic-side prediction is \(a_0(0)=(c/2)\sqrt{G\rho_\mathrm{DE,0}}\) with
\(\rho_\mathrm{DE,0}=\Omega_\mathrm{DE}\,3H_0^2/(8\pi G)\):

{\def\LTcaptype{none} % do not increment counter
\begin{longtable}[]{@{}
  >{\raggedright\arraybackslash}p{(\linewidth - 6\tabcolsep) * \real{0.2500}}
  >{\raggedright\arraybackslash}p{(\linewidth - 6\tabcolsep) * \real{0.2500}}
  >{\raggedright\arraybackslash}p{(\linewidth - 6\tabcolsep) * \real{0.2500}}
  >{\raggedright\arraybackslash}p{(\linewidth - 6\tabcolsep) * \real{0.2500}}@{}}
\toprule\noalign{}
\begin{minipage}[b]{\linewidth}\raggedright
footing
\end{minipage} & \begin{minipage}[b]{\linewidth}\raggedright
cosmic \(a_0(0)\)
\end{minipage} & \begin{minipage}[b]{\linewidth}\raggedright
\(a_{0,\rm SPARC}/a_{0,\rm cosmic}\)
\end{minipage} & \begin{minipage}[b]{\linewidth}\raggedright
tension
\end{minipage} \\
\midrule\noalign{}
\endhead
\bottomrule\noalign{}
\endlastfoot
Planck (\(\Omega_\mathrm{DE}=0.685,\,H_0=67.4\)) & \(9.36\times10^{-11}\) &
\(1.26\) & \(1.44\sigma\) \\
SH0ES-\(H_0\) (\(\Omega_\mathrm{DE}=0.685,\,H_0=73.0\)) &
\(1.01\times10^{-10}\) & \(1.17\) & \(0.95\sigma\) \\
\end{longtable}
}

The tie \textbf{holds at \(\sim1\sigma\)} on the canonical
(Planck-\(H_0\)) footing and slightly better (\(0.95\sigma\)) at the
local \(H_0\). This is by construction --- the definitional anchor ---
and holds. It re-states that the framework enters the high-\(z\) test on
solid \(z=0\) footing, and carries no evidential weight beyond that
internal consistency.

\subsubsection{3.2 The cosmology side: DESI DR2
propagated}\label{the-cosmology-side-desi-dr2-propagated}

On the cosmic side there is no free dark-energy parameter to fit --- the
framework consumes whatever \(\rho_\mathrm{DE}(z)\) the data prefer. We
propagate the DESI DR2 (2025, arXiv:2503.14738) \(w_0w_a\)CDM posterior
(BAO + CMB + SNe; representative marginals with \(w_0\)--\(w_a\)
correlation \(\approx-0.87\); \(\Lambda\)CDM excluded at
\(2.8\sigma\)/Pantheon+, \(4.2\sigma\)/DESY5, \(3.8\sigma\)/Union3)
through \(a_0(z)/a_0(0)=\sqrt{\rho_\mathrm{DE}(z)/\rho_\mathrm{DE,0}}\)
(Appendix A, \texttt{desi\_posterior\_a0z.py}):

{\def\LTcaptype{none} % do not increment counter
\begin{longtable}[]{@{}
  >{\raggedright\arraybackslash}p{(\linewidth - 10\tabcolsep) * \real{0.1667}}
  >{\raggedright\arraybackslash}p{(\linewidth - 10\tabcolsep) * \real{0.1667}}
  >{\raggedright\arraybackslash}p{(\linewidth - 10\tabcolsep) * \real{0.1667}}
  >{\raggedright\arraybackslash}p{(\linewidth - 10\tabcolsep) * \real{0.1667}}
  >{\raggedright\arraybackslash}p{(\linewidth - 10\tabcolsep) * \real{0.1667}}
  >{\raggedright\arraybackslash}p{(\linewidth - 10\tabcolsep) * \real{0.1667}}@{}}
\toprule\noalign{}
\begin{minipage}[b]{\linewidth}\raggedright
SNe combination
\end{minipage} & \begin{minipage}[b]{\linewidth}\raggedright
\(a_0(3)/a_0(0)\)
\end{minipage} & \begin{minipage}[b]{\linewidth}\raggedright
\([16\%,84\%]\)
\end{minipage} & \begin{minipage}[b]{\linewidth}\raggedright
\(P(a_0\text{ declines})\)
\end{minipage} & \begin{minipage}[b]{\linewidth}\raggedright
decline signif.
\end{minipage} & \begin{minipage}[b]{\linewidth}\raggedright
in \(0.60\)--\(0.75\)
\end{minipage} \\
\midrule\noalign{}
\endhead
\bottomrule\noalign{}
\endlastfoot
DESI+CMB+Pantheon+ & \(0.775\) & \([0.683,\,0.878]\) & \(97.9\%\) &
\(2.0\sigma\) & \(38\%\) \\
DESI+CMB+DESY5 & \(0.737\) & \([0.651,\,0.833]\) & \(99.3\%\) &
\(2.5\sigma\) & \(51\%\) \\
DESI+CMB+Union3 & \(0.707\) & \([0.599,\,0.831]\) & \(98.3\%\) &
\(2.1\sigma\) & \(48\%\) \\
\end{longtable}
}

Across the three combinations the median is \(a_0(3)/a_0(0)\simeq0.74\)
with a \(\sim2.2\sigma\) decline, and all three posteriors overlap the
framework's \(0.60\)--\(0.75\) band. This is a genuine, \emph{live}
cosmology-side hint: \textbf{if} the DESI evolving-dark-energy signal is
real, the induced \(a_0(z)\) decline falls squarely inside the
framework's \(0.60\)--\(0.75\) band. It is not, however, an independent
test --- it is a re-expression of the DESI \(w(z)\) result through the
framework relation, the \emph{same} DESI information seen twice, and it
therefore \textbf{cannot be added to the galaxy-side \(\chi^2\)} as if
it were a second measurement. It dissolves entirely if the DESI signal
regresses to \(w=-1\). It sharpens the \emph{target} for the galaxy
side; it does not by itself confront it.

\begin{center}\rule{0.5\linewidth}{0.5pt}\end{center}

\subsection{\texorpdfstring{4. The galaxy-side \(a_0(z)\) from the
high-\(z\)
BTFR}{4. The galaxy-side a\_0(z) from the high-z BTFR}}\label{the-galaxy-side-a_0z-from-the-high-z-btfr}

\subsubsection{4.1 The estimator}\label{the-estimator}

On the galaxy side, \(a_0(z)\) is read from the BTFR zero-point at each
redshift. In the deep-MOND regime \(V^4=a_0\,GM_\mathrm{bar}\) is exact
(the fitted slope is forced to \(4\)), so at fixed circular velocity

\[\frac{a_0(z)}{a_0(0)}=\frac{(V_z/V_0)^4}{M_{\mathrm{bar},z}/M_{\mathrm{bar},0}},\qquad\text{equivalently}\qquad \Delta\log a_0=-\,\Delta\log M_\mathrm{bar}\big|_{\text{fixed }V}.\]

Reading \(a_0(z)\) is therefore the same operation as reading the BTFR
zero-point. The estimator is valid \textbf{only} where the fitted slope
is genuinely \(4\), i.e.~for objects deep in the low-acceleration regime
\(g\ll a_0\).

\subsubsection{\texorpdfstring{4.2 The one clean high-\(z\) datum: the
Big Wheel at
\(z=3.25\)}{4.2 The one clean high-z datum: the Big Wheel at z=3.25}}\label{the-one-clean-high-z-datum-the-big-wheel-at-z3.25}

A single clean deep-MOND high-\(z\) rotator is currently available ---
the Big Wheel (Wang et al.~2025, \emph{Nature Astronomy} \textbf{9},
710; \(z=3.2514\); arXiv:2409.17956), a large, low-acceleration disk
whose outer rotation curve satisfies the deep-MOND condition.
Monte-Carlo over its velocity and baryonic-mass systematics
(\(V_\mathrm{rot}=280\pm31\,{\rm km\,s^{-1}}\),
\(\sigma_\mathrm{int}=61\,{\rm km\,s^{-1}}\) with asymmetric-drift
correction; \(M_\star=1.7\pm0.7\times10^{11}M_\odot\), dynamically
capped; \(M_\mathrm{gas}=1.8\pm0.9\times10^{11}M_\odot\), \(\times1.36\)
He) gives
\(a_{0,\rm eff}=1.54\,(+1.10/-0.61)\times10^{-10}\,{\rm m\,s^{-2}}\),
i.e.

{\def\LTcaptype{none} % do not increment counter
\begin{longtable}[]{@{}
  >{\raggedright\arraybackslash}p{(\linewidth - 2\tabcolsep) * \real{0.5000}}
  >{\raggedright\arraybackslash}p{(\linewidth - 2\tabcolsep) * \real{0.5000}}@{}}
\toprule\noalign{}
\begin{minipage}[b]{\linewidth}\raggedright
footing
\end{minipage} & \begin{minipage}[b]{\linewidth}\raggedright
\(a_0(3.25)/a_0(0)\)
\end{minipage} \\
\midrule\noalign{}
\endhead
\bottomrule\noalign{}
\endlastfoot
SPARC \(1.181\times10^{-10}\) & \(1.31\,(+0.93/-0.52)\) \\
canonical \(0.936\times10^{-10}\) (\(\rho_\mathrm{DE}/cH_\Lambda\)) &
\(1.65\,(+1.17/-0.65)\) \\
ALT \(1.131\times10^{-10}\) (\(\rho_\mathrm{tot}/cH_0\)) &
\(1.37\,(+0.97/-0.54)\) \\
\end{longtable}
}

On \textbf{every} footing the ratio is consistent with a constant
(\(1.0\)) and with the DESI decline (\(0.68\)), while it
\textbf{disfavors the rising \(cH(z)\) evolution model (\(\sim5\times\))
at \(\sim2\sigma\)} (banked posterior probability \(\sim1\)--\(3\%\);
the exclusion is robust because reaching \(5\times\) would require
\(M_\mathrm{bar}\) wrong by \(\sim0.7\) dex, far outside any systematic).
This is a real, if narrow, result: one clean object already kills the
rising-evolution branch.

What it cannot do is separate the \(0.70\) decline from flat. The
\(M_\mathrm{bar}\) systematic on a single object (\(\approx0.3\) dex in
\(M_\mathrm{bar}\), propagating to \(\approx0.23\) dex in the combined
\(a_0\) error) is itself larger than the \(0.164\)-dex flat-vs-decline
signal at \(z=3.25\). Worse, that \(0.226\)-dex error is a \emph{floor}:
it holds the asymmetric-drift coefficient fixed at \(\alpha=3.4\) rather
than Monte-Carloing it, whereas varying \(\alpha\) moves the ratio from
\(0.95\) (\(\alpha=0\)) to \(1.27\) (\(\alpha=3.4\)) --- a further
\(\sim0.1\)--\(0.15\) dex --- so the true error is somewhat larger and
the datum even less constraining (the omitted term also pushes \(a_0\)
\emph{up}, i.e.~anti-framework, so it does not rescue a ``win''). We use
the anti-framework choices at every fork: the dynamically capped stellar
mass, which pushes the ratio \emph{up} (the Monte-Carlo median is
\(1.31\); the capped-mass point value alone is \(1.27\); the full-SED
stellar mass would instead place it at \(0.86\), right on the DESI
decline), and a high asymmetric-drift correction --- and the datum still
lands consistent with both tracks.

One further systematic works in the framework's favour and is \emph{not}
applied to the central value, so the quoted ratio is conservative. The
estimator \(a_{0,\rm eff}\equiv V^4/(GM_\mathrm{bar})\) equals \(a_0\)
\emph{only} in the deep-MOND limit \(g_\mathrm{bar}\ll a_0\). On the
framework's own interpolation
\(g_\mathrm{obs}^2=g_\mathrm{bar}^2+g_\mathrm{bar}a_0\) one has exactly
\(a_{0,\rm eff}=g_\mathrm{obs}^2/g_\mathrm{bar}=g_\mathrm{bar}+a_0\), so the
estimator \emph{overstates} \(a_0\) by \(g_\mathrm{bar}\) whenever the
rotation-curve point is not deeply in the MOND regime. For the Big Wheel
(\(M_\mathrm{bar}\approx4.1\times10^{11}M_\odot\), outer point at
\(R\sim30\)--\(45\) kpc) \(g_\mathrm{bar}\approx0.3\)--\(0.5\,a_0\), so
\(a_{0,\rm eff}\) is biased high by \(\sim30\)--\(50\%\): subtracting
\(g_\mathrm{bar}\) moves the datum from
\(a_{0,\rm eff}=1.54\times10^{-10}\) (ratio \(1.31\)) toward
\(a_0\approx1.1\)--\(1.2\times10^{-10}\) (ratio \(\approx1.0\)),
i.e.~\emph{closer} to flat and to the decline, and \emph{further} from
the rise. The ``clean deep-MOND'' label is therefore optimistic for this
single transitional object --- the correct treatment uses the full
\(\nu\)-kernel at the measured \(g_\mathrm{bar}/a_0\), not the deep-MOND
asymptote --- but the direction of the correction only reinforces the
``consistent-with-both, rise-excluded'' reading, and the actual outer-RC
radius (hence \(g_\mathrm{bar}/a_0\)) should be pinned before this datum is
used quantitatively.

\subsubsection{\texorpdfstring{4.3 Why intermediate-\(z\) is not a clean
\(a_0\)
probe}{4.3 Why intermediate-z is not a clean a\_0 probe}}\label{why-intermediate-z-is-not-a-clean-a_0-probe}

The temptation is to add the many \(z\sim0.5\)--\(2.3\) rotating disks
(e.g.~KMOS\(^{3\rm D}\)). They do not help, for a physical reason: these
massive rotators sit at \(g\gtrsim a_0\), \textbf{not} deep-MOND, so
their fitted BTFR slopes are \(3.0\)--\(3.85\), not \(4\), and the
zero-point is then degenerate with the slope, the pivot, and
baryon/dark-matter-fraction evolution. The inferred ``\(a_0\)
evolution'' \emph{flips sign with the analysis choice}:

\begin{itemize}
\tightlist
\item
  Übler et al.~(2017, KMOS\(^{3\rm D}\), fixed local slope \(3.75\))
  \(\to\) zero-point \(\sim0.3\)--\(0.4\) dex below local \(\to\) naive
  map reads \(a_0\) \textbf{rising} \(\sim2\)--\(3\times\);
\item
  Sharma et al.~(2024, free slope \(3.21\)) \(\to\) zero-point above
  local at their pivot \(\to\) \textbf{opposite} sign;
\item
  Di Teodoro et al.~(2016) and Tiley et al.~(2019, \(V/\sigma>3\))
  \(\to\) \textbf{null}, no significant zero-point evolution.
\end{itemize}

The scatter across analyses (a factor \(\sim2\)--\(3\), in \emph{either}
direction) dwarfs the \(0.70\)-vs-\(1.0\) signal. These points are not
valid \(a_0\) readouts; where we include the \(z\approx2.3\) point in
the joint fit we place it at its fixed-slope value \(1.86\) (rising ---
the \emph{wrong} sign for the framework decline), so it counts against,
not for, the framework.

\subsubsection{\texorpdfstring{4.4 The \(a_0=M_\mathrm{bar}\) degeneracy
and the mild distance
dependence}{4.4 The a\_0=M\_\{\textbackslash rm bar\} degeneracy and the mild distance dependence}}\label{the-a_0m_rm-bar-degeneracy-and-the-mild-distance-dependence}

Because \(\Delta\log a_0=-\Delta\log M_\mathrm{bar}\) at fixed \(V\), a
mis-estimate of the baryonic mass propagates \(1{:}1\) into \(a_0\): the
two are algebraically degenerate for this estimator. Rising
molecular-gas fractions (\(\sim10\%\) at \(z=0\) to \(\sim50\%\) at
\(z=2\)), the CO-to-H\(_2\) conversion \(\alpha_\mathrm{CO}\), dust, and
IMF evolution all move the BTFR zero-point and are indistinguishable
from an \(a_0\) shift with the BTFR / point-mass method. This is the
crux limitation of §5. The mild distance--cosmology dependence is
separate and quantified: \(M_\mathrm{bar}\propto D^2\), and the ratio
\(D(\text{DESI})/D(\text{flat})\) shifts \(M_\mathrm{bar}\) by \(\le0.015\)
dex at \(z=1\) (smaller at higher \(z\)) --- roughly \(10\times\) below
the signal and \(15\times\) below the Big-Wheel error. A wrong
background moves all points together and cannot manufacture a decline,
so the test is \textbf{not circular in the fatal sense}: testing whether
\(a_0\) tracks \(\rho_\mathrm{DE}\) is not the same as assuming the DE
evolution.

\subsubsection{4.5 Does the galaxy data show the
decline?}\label{does-the-galaxy-data-show-the-decline}

No --- neither way. Assembling the honest galaxy points (Appendix A,
\texttt{confront.py}):

{\def\LTcaptype{none} % do not increment counter
\begin{longtable}[]{@{}
  >{\raggedright\arraybackslash}p{(\linewidth - 8\tabcolsep) * \real{0.2000}}
  >{\raggedright\arraybackslash}p{(\linewidth - 8\tabcolsep) * \real{0.2000}}
  >{\raggedright\arraybackslash}p{(\linewidth - 8\tabcolsep) * \real{0.2000}}
  >{\raggedright\arraybackslash}p{(\linewidth - 8\tabcolsep) * \real{0.2000}}
  >{\raggedright\arraybackslash}p{(\linewidth - 8\tabcolsep) * \real{0.2000}}@{}}
\toprule\noalign{}
\begin{minipage}[b]{\linewidth}\raggedright
set
\end{minipage} & \begin{minipage}[b]{\linewidth}\raggedright
\(\chi^2(\text{flat})\)
\end{minipage} & \begin{minipage}[b]{\linewidth}\raggedright
\(\chi^2(\text{DESI-decl})\)
\end{minipage} & \begin{minipage}[b]{\linewidth}\raggedright
\(\Delta\chi^2\)
\end{minipage} & \begin{minipage}[b]{\linewidth}\raggedright
lean
\end{minipage} \\
\midrule\noalign{}
\endhead
\bottomrule\noalign{}
\endlastfoot
all 4 points & \(1.08\) & \(3.05\) & \(-1.97\) & flat, weak \\
clean deep-MOND only (\(z=0\) + Big Wheel) & \(0.27\) & \(1.55\) &
\(-1.28\) & flat, weak \\
\end{longtable}
}

The clean-only row (\(z=0\) + Big Wheel) is the sole formally meaningful
statistic; the ``all 4 points'' row is \textbf{illustrative only},
because the intermediate-\(z\) points are declared invalid \(a_0\)
readouts (§4.3, slope \(\ne4\), sign-contested) and the \(z=1\) entry is
itself the median of a hand-picked literature straddle --- quantities
that cannot legitimately enter a likelihood. Both rows agree in
direction: the data mildly prefer flat,
\(|\Delta\chi^2|\approx1.3\)--\(2.0\) --- not significant, and what lean
exists in the intermediate-\(z\) points is toward \emph{rising}, the
wrong sign for the framework decline. The distinctive \(0.60\)--\(0.75\)
decline is not detected; the data are consistent with both flat and the
decline. The one branch the data exclude is the rising \(cH(z)\)
evolution model (\(\sim2\sigma\), via the Big Wheel).

\begin{center}\rule{0.5\linewidth}{0.5pt}\end{center}

\subsection{\texorpdfstring{5. The key finding: \(M_\mathrm{bar}\)
calibration, not rotator count, is
decisive}{5. The key finding: M\_\{\textbackslash rm bar\} calibration, not rotator count, is decisive}}\label{the-key-finding-m_rm-bar-calibration-not-rotator-count-is-decisive}

The decisive question for any future test is how the significance
scales. The separability statistic --- the maximum sigma at which the
\emph{current} errors could ever tell flat from the DESI decline --- is

\[S=\sqrt{\sum_i\left[\frac{\log a_0^\mathrm{flat}(z_i)-\log a_0^\mathrm{DESI}(z_i)}{\sigma_i}\right]^2}.\]

The current data give \(S=0.73\sigma\) (Big Wheel alone), \(0.73\sigma\)
(clean pair \(z=0+3.25\)), and \(0.80\sigma\) (all four points):
\textbf{underpowered}, consistent with both tracks. Every galaxy point
sits \(\le1.2\sigma\) from \emph{both} the flat and the declining track.
So more data are needed --- but the crucial point is \emph{what kind}.

The naive forecast assumes independent per-object errors that beat down
as \(1/\sqrt{N}\). The target signal is \(-\log_{10}(0.70)=0.155\) dex,
requiring \(\sigma_\mathrm{mean}=0.052\) dex for \(3\sigma\), hence
\(N\approx4\) rotators at SPARC-like precision (\(0.10\) dex) or
\(N\approx24\) at realistic high-\(z\) per-object precision (\(0.25\)
dex). \textbf{This is wrong}, because the dominant \(M_\mathrm{bar}\) error
--- \(\alpha_\mathrm{CO}\), IMF, and dust prescriptions applied
\emph{identically} to a sample --- is largely common-mode, not
independent, and a common-mode error does not average down with \(N\).
With \(\sigma_\mathrm{mean}^2=\sigma_\mathrm{ind}^2/N+\sigma_\mathrm{cm}^2\) and
a correlated \(M_\mathrm{bar}\) floor \(\sigma_\mathrm{cm}\):

{\def\LTcaptype{none} % do not increment counter
\begin{longtable}[]{@{}ll@{}}
\toprule\noalign{}
\(M_\mathrm{bar}\) common-mode floor & \(N\) for \(3\sigma\) \\
\midrule\noalign{}
\endhead
\bottomrule\noalign{}
\endlastfoot
\(0.00\) dex & \(24\) \\
\(0.05\) dex & \(\sim377\) \\
\(\ge0.08\) dex & \textbf{never reaches \(3\sigma\) at any \(N\)} \\
\end{longtable}
}

Because realistic high-\(z\) \(M_\mathrm{bar}\) calibration floors are
\(\sim0.08\)--\(0.15\) dex (\(\alpha_\mathrm{CO}\) alone is a
\(0.1\)--\(0.2\) dex systematic), the test is decisive \textbf{only if
the common-mode \(M_\mathrm{bar}\) systematic can be driven below
\(\sim0.05\) dex} --- a calibration problem, not a counting problem.
This is the paper's core methodological contribution: \textbf{the
decisive observational lever is a uniformly calibrated \(M_\mathrm{bar}\)
ladder to \(\lesssim0.05\) dex common-mode, not the number of rotators.}
A uniformly calibrated \(M_\mathrm{bar}\) scale is worth more than dozens
of noisy objects. (Both footings enter this identically --- the forecast
is expressed in ratios and dex, so the \(20.9\%\) footing choice
cancels; it re-enters only the Big-Wheel absolute of §4.2.)

\begin{center}\rule{0.5\linewidth}{0.5pt}\end{center}

\subsection{\texorpdfstring{6. Forecast and the \(a_0\)-separable
observable}{6. Forecast and the a\_0-separable observable}}\label{forecast-and-the-a_0-separable-observable}

To reach \(3\sigma\) on the \(0.60\)--\(0.75\) decline versus flat
therefore requires (i) a common-mode \(M_\mathrm{bar}\) calibration
\(\lesssim0.05\) dex and (ii) a sample of \(\sim20\)--\(40\) clean
deep-MOND \(z\sim2\)--\(3\) rotators, Big-Wheel-like, with gas-traced
outer rotation curves (JWST/ALMA) or individually resolved deep-MOND
curves (ELT/HARMONI). For fixed samples the per-object precision
required is \(\le0.16\) dex (\(N=10\)), \(\le0.23\) dex (\(N=20\)),
\(\le0.33\) dex (\(N=40\)) --- all under the assumption that these
errors are independent, which returns the argument to the common-mode
floor above. Clean deep-MOND disks like the Big Wheel are rare (the
intermediate-\(z\) KMOS\(^{3\rm D}\) rotators are \(g\gtrsim a_0\),
§4.3), so assembling \(20\)--\(40\) of them is itself an observational
stretch.

The deeper methodological escape is to stop using the BTFR zero-point at
all. Because \(a_0\) and \(M_\mathrm{bar}\) are \(1{:}1\) degenerate in the
BTFR, no amount of BTFR data isolates \(a_0\) from baryonic-mass
evolution. The \textbf{\(a_0\)-separable} observable is the \emph{shape}
of the rotation curve --- the radial-acceleration-relation transition
radius / the full RAR --- which depends on \(a_0\) at fixed
\(M_\mathrm{bar}\) and therefore breaks the degeneracy that the zero-point
cannot. A high-\(z\) RAR-transition measurement (the acceleration at
which the curve departs from the Newtonian-baryonic expectation) reads
\(a_0\) directly, without hostage to the gas-mass normalization. This is
the observable to target: not more BTFR zero-points, but resolved
high-\(z\) rotation-curve shapes.

\begin{center}\rule{0.5\linewidth}{0.5pt}\end{center}

\subsection{7. Discussion}\label{discussion}

The cross-scale link is the substance here. \(\Lambda\)CDM contains no
relation forcing a galaxy's internal acceleration scale to equal, or to
co-evolve with, the cosmic dark-energy density; the framework forces
both. That makes the joint \(z\)-tracking of galaxy kinematics and
cosmic distances a genuine, falsifiable prediction that is not available
to \(\Lambda\)CDM --- and it is non-circular precisely because the two
datasets never share inputs. The present status is asymmetric between
the two scales. On the \textbf{cosmology side} the test is live: the
DESI DR2 evolving-dark-energy posterior, propagated through the
framework relation, produces an \(a_0(3)/a_0(0)\simeq0.74\) decline (at
\(\sim2.2\sigma\)) that falls inside the framework's \(0.60\)--\(0.75\)
band --- but this is a consequence of the framework relation plus the
DESI \(w(z)\), not an independent measurement, and it stands or falls
with the DESI signal. On the \textbf{galaxy side} the test is a hostage
to \(M_\mathrm{bar}\): the distinctive decline is neither detected nor
excluded (\(S\simeq0.8\sigma\)), the intermediate-\(z\) BTFR is
systematics-dominated and sign-contested, and the one clean high-\(z\)
datum can only exclude the rising ALT footing. The honest reading is
that the framework's forced bridge is currently underpowered --- not
passed, not failed --- and that the binding constraint is
\(M_\mathrm{bar}\) calibration rather than sample size.

Two caveats bound the interpretation. First, the \(a_0\) magnitude
inherits the posited \(Z\); only the ratio/tracking is under test, so a
detection of the decline would test the \emph{relation}, not the
\emph{value} of \(a_0\) or \(Z\). Second, the \(z=0\) tie is
definitional and carries no independent weight; the evidential content
is entirely in the slope. The framework is honest both ways: at every
fork in the Big-Wheel analysis the anti-framework option is taken, and
the datum still lands consistent with both tracks; the adverse rising
intermediate-\(z\) point is retained in the fit; and the forecast is
corrected downward by the common-mode floor.

\begin{center}\rule{0.5\linewidth}{0.5pt}\end{center}

\subsection{8. Conclusion}\label{conclusion}

We have set out a non-circular cross-scale test of de Sitter--Unruh
modified inertia: the framework replaces \(\Lambda\)CDM's free
dark-energy fit with the galaxy acceleration scale,
\(H^2=\tfrac{8\pi G}{3}\rho_m+Z^2a_0^2(z)/c^2\), forcing
\(a_0(z)/a_0(0)=\sqrt{\rho_\mathrm{DE}(z)/\rho_\mathrm{DE,0}}\) --- a bridge
between galaxy kinematics and cosmic distances that \(\Lambda\)CDM
structurally lacks. The \(z=0\) tie holds at \(\sim1\sigma\) (SPARC
\(a_0(0)=1.181\times10^{-10}\) vs cosmic
\((c/2)\sqrt{G\rho_\mathrm{DE,0}}\); ratio \(1.26\) Planck-\(H_0\),
\(1.17\) SH0ES-\(H_0\)), the definitional anchor. The cosmology side is
live: the DESI DR2 posterior propagated through the relation yields
\(a_0(3)/a_0(0)\simeq0.74\) at \(\sim2.2\sigma\) --- a consequence of
the framework tie plus the DESI \(w(z)\), not an independent detection,
and it dissolves if \(w\to-1\). The galaxy side is underpowered: the
distinctive \(0.60\)--\(0.75\) decline is neither detected nor excluded
(\(S\simeq0.8\sigma\)); one clean \(z=3.25\) rotator excludes only the
ALT-\(cH_0\) rise (\(\sim2\sigma\)). The central finding is
methodological and redirects the observational strategy: in the BTFR
estimator \(a_0\) is \(1{:}1\) degenerate with \(M_\mathrm{bar}\), the
dominant systematic is common-mode, and with a \(\gtrsim0.08\)-dex
correlated \(M_\mathrm{bar}\) floor the test never reaches \(3\sigma\) at
any rotator count. The decisive lever is \(M_\mathrm{bar}\) calibration to
\(\lesssim0.05\) dex common-mode, and the \(a_0\)-separable observable
is the full rotation-curve / RAR transition, not the BTFR zero-point.
The value of \(a_0\) and the constant \(Z\) are posited, not derived;
both footings are carried throughout; and the test is reported honestly
as a non-detection with a live cosmology-side hint and a strategic
redirection of the future program.

\begin{center}\rule{0.5\linewidth}{0.5pt}\end{center}

\subsection{Appendix A: The committed, exit-0 verification
scripts}\label{appendix-a-the-committed-exit-0-verification-scripts}

All numbers above are reproduced by three committed scripts (numpy/scipy
only, exit 0, both footings), re-run 2026-07-18; the frozen public
repository is left read-only and all outputs are written to
\texttt{prep\_2026/a0z\_crossscale/}.

\textbf{\texttt{galaxy\_a0z.py}} --- the \(\Lambda\)-blind galaxy-side
\(a_0(z)\) from the high-\(z\) BTFR zero-point. Assembles the \(z=0\)
SPARC anchor (\(1.181\times10^{-10}\pm16\%\)), the sign-contested
intermediate-\(z\) compilation (Übler+2017, Sharma+2024, Di
Teodoro+2016, Tiley+2019) with the deep-MOND-validity flag, and the
\(z=3.25\) Big-Wheel Monte-Carlo
(\(a_{0,\rm eff}=1.54\,(+1.10/-0.61)\times10^{-10}\)). Reports ratios so
the footing cancels, and re-enters only the Big-Wheel absolute
(SPARC/canonical/ALT \(=1.31/1.65/1.37\)). Prints the forecast
\(N\approx4\) (clean) / \(N\approx24\) (realistic) and writes
\texttt{galaxy\_a0z.png}.

\textbf{\texttt{confront.py}} --- confronts the galaxy \(a_0(z)\)
against the flat-\(\Lambda\) and DESI-DR2 evolving cosmic tracks.
Produces the joint \(\chi^2\) (all-4: \(1.08\) vs \(3.05\); clean:
\(0.27\) vs \(1.55\)), the separability (\(S=0.73/0.73/0.80\sigma\)),
the \(z=0\) tie (\(1.26/1.44\sigma\) Planck, \(1.17/0.95\sigma\) SH0ES),
the both-footings Big-Wheel absolute, and the forecast table --- the
fixed-\(N\) precision requirements \textbf{and} the common-mode
\(M_\mathrm{bar}\) floor of §5 (\(\sigma_\mathrm{cm}=0\to N\!\approx\!24\);
\(0.05\to N\!\approx\!377\); \(\ge0.08\to\) never reaches \(3\sigma\) at
any \(N\)). Bottom line: \(S(\text{all})=0.80\sigma\Rightarrow\)
underpowered, consistent with both.

\textbf{\texttt{desi\_posterior\_a0z.py}} --- propagates the DESI DR2
\(w_0w_a\)CDM posterior (three SNe combinations, correlated
\(w_0\)--\(w_a\) via a seeded Cholesky draw) through
\(a_0(z)/a_0(0)=\sqrt{\rho_\mathrm{DE}(z)/\rho_\mathrm{DE,0}}\), giving
\(a_0(3)/a_0(0)=0.775\) (Pantheon+, \(2.0\sigma\)), \(0.737\) (DESY5,
\(2.5\sigma\)), \(0.707\) (Union3, \(2.1\sigma\)); combined median
\(\simeq0.74\), decline significance \(\sim2.2\sigma\), all overlapping
the framework \(0.60\)--\(0.75\) band. The ratio is footing-, \(H_0\)-,
and \(M_B\)-independent (canonical \(\rho_\mathrm{DE}\) footing by
construction; \(Z\) cancels).

Supporting reconstructions from an independent lane
(\texttt{prep\_2026/a0z\_from\_sne/}, \texttt{crossscale.py}) confirm
that Type Ia supernovae \textbf{alone} cannot measure \(a_0(z=3)\): the
dark-energy residual \(E^2-\Omega_m(1+z)^3\) is a difference of large
numbers, and the model-independent Gaussian-process reconstruction
(Seikel, Clarkson \& Smith 2012) is sign-pinned only to
\(z\approx0.40\). The robust SNe deliverable is \(a_0(0)\)
(\(\simeq9.4\times10^{-11}\) at \(H_0=67.4\),
\(\simeq1.01\times10^{-10}\) at \(H_0=73.0\)), agreeing with SPARC at
\(\sim1\sigma\); the \(z=3\) decline numbers come from the DESI/model
references, not a SNe reconstruction --- which is exactly why the
two-dataset (galaxy \(\times\) cosmic-distance) construction, and not
SNe alone, is the test.

\begin{center}\rule{0.5\linewidth}{0.5pt}\end{center}

\subsection{References}\label{references}

Brout, D., et al.~2022, \emph{ApJ} \textbf{938}, 110 (Pantheon+).

DESI Collaboration 2025, arXiv:2503.14738 (\(w_0w_a\)CDM; \(\Lambda\)CDM
excluded at \(2.8/3.8/4.2\sigma\)).

Di Teodoro, E. M., Fraternali, F., \& Miller, S. H. 2016, \emph{A\&A}
\textbf{594}, A77.

Lelli, F., McGaugh, S. S., \& Schombert, J. M. 2016, \emph{AJ}
\textbf{152}, 157 (SPARC).

McGaugh, S. S. 2005, \emph{ApJ} \textbf{632}, 859; McGaugh, S. S. 2012,
\emph{AJ} \textbf{143}, 40 (BTFR).

Milgrom, M. 1983, \emph{ApJ} \textbf{270}, 365; Milgrom, M. 1999,
\emph{Phys. Lett. A} \textbf{253}, 273 (the \(\nu\)-kernel, Eq. 9;
astro-ph/9805346).

Seikel, M., Clarkson, C., \& Smith, M. 2012, \emph{JCAP} \textbf{06},
036 (Gaussian-process \(H(z)\)).

Sharma, G., Upadhyaya, V., Salucci, P., \& Desai, S. 2024, \emph{A\&A}
\textbf{690} (Tully--Fisher at \(0.6\le z\le2.5\); free bTFR slope
\(3.21\pm0.28\); arXiv:2406.08934).

Tiley, A. L., et al.~2019, \emph{MNRAS} \textbf{482}, 2166 (KROSS--SAMI
Tully--Fisher across 8 Gyr since \(z\approx1\)).

Übler, H., et al.~2017, \emph{ApJ} \textbf{842}, 121 (KMOS\(^{3\rm D}\)
bTFR).

Wang, W., et al.~2025, \emph{Nature Astronomy} \textbf{9}, 710 (the
``Big Wheel'' giant disk at \(z=3.2514\); arXiv:2409.17956).

Zimmerman, C. 2026, \emph{Reading the Cosmological Constant from
Dwarf-Galaxy Rotation Curves: The \(a_0\)-Line} (self-cite), DOI
10.5281/zenodo.21419735.

Zimmerman, C. 2026, \emph{MI Field Theory Results 2026} (self-cite), DOI
10.5281/zenodo.21403470.

\end{document}
